\documentclass[]{foi} 
\usepackage[utf8]{inputenc}
\usepackage[T1]{fontenc}
\usepackage[croatian]{babel}
\DeclareUnicodeCharacter{2011}{-}
\usepackage{lipsum}
\usepackage{float}

\vrstaRada{\seminar}

\title{API za rad s različitim formatima datoteka Apache POI}
\predmet{Napredne Web tehnologije i servisi}

\author{Zlatko Pračić} 
\spolStudenta{\musko} 

\mentor{Dragutin Kermek}
\spolMentora{\musko} 
\titulaProfesora{Prof. dr. sc.}

\godina{2026}
\mjesec{kolovoz}

\indeks{0016145097}

\smjer{Baze podataka i baze znanja}



\sazetak{Opsega od 100 do 300 riječi. Sažetak upućuje na temu rada, ukratko se iznosi čime se rad bavi, teorijsko-metodološka polazišta, glavne teze i smjer rada te zaključci.}

\kljucneRijeci{riječ; riječ; riječ; Obuhvaća $7\pm2$ ključna pojma koji su glavni predmet rasprave u radu.}

\begin{document}

\maketitle

\tableofcontents

\makeatletter \def\@dotsep{4.5} \makeatother
\pagestyle{plain}

\chapter{Uvod}
U suvremenom okruženju razmjena podataka putem Microsoft Office dokumenata svakodnevna je pojava. Excel tablice, Word dokumenti i PowerPoint prezentacije postali su de facto standard za izvještavanje, analizu podataka i komunikaciju. Međutim, ručna obrada takvih dokumenata u sustavima koji zahtijevaju automatizaciju, poput sustava za upravljanje sadržajem, poslovnih aplikacija ili alata za ekstrakciju podataka, nije praktična ni skalabilna.

Apache POI projekt, koji razvija i održava Apache Software Foundation, nudi rješenje za ovaj problem u obliku besplatne Java biblioteke otvorenog koda. Biblioteka programerima omogućuje čitanje, pisanje i izmjenu Microsoft Office datoteka isključivo korištenjem Java koda, bez potrebe za instaliranim Microsoft Officeom. Time se postiže potpuna platformska neovisnost, što je posebno vrijedno u enterprise okruženjima koja se oslanjaju na Java ekosustav.

Motivacija za odabir Apache POI-a kao teme ovog rada leži upravo u njegovoj rasprostranjenosti i praktičnoj primjenjivosti. Biblioteka je jedna od najkorištenijih Java biblioteka za rad s Office formatima i aktivno se razvija već više od dva desetljeća. Podržava kako starije binarne formate temeljene na OLE2 standardu (XLS, DOC, PPT), tako i moderne XML-bazirane OOXML formate uvedene s Microsoft Officeom 2007 (XLSX, DOCX, PPTX).

Cilj ovog rada je sustavno prikazati arhitekturu i ključne komponente Apache POI projekta te demonstrirati njegovu primjenu na konkretnim primjerima. Rad je strukturiran na sljedeći način: drugo poglavlje donosi pregled povijesti projekta i njegovih tehničkih komponenti, treće poglavlje opisuje rad s pojedinim formatima datoteka, a četvrto poglavlje prikazuje praktične primjere generiranja Excel izvještaja i obrade Word dokumenta. Na kraju se iznosi zaključak s osvrtom na prednosti, ograničenja i smjernice za primjenu biblioteke.

\chapter{Pregled Apache POI projekta}

\section{Povijest i razvoj}
Od Jakarta POI projekta do današnjeg Apache POI-a.

\section{Arhitektura i komponente}
Pregled modula (HSSF, XSSF, SXSSF, HWPF, XWPF, HSLF, XSLF) i razlika 
između binarnih OLE2 i modernih OOXML implementacija.

\chapter{Rad s formatima datoteka}

\section{Excel datoteke}
Čitanje i pisanje .xls (HSSF) i .xlsx (XSSF) formata, formatiranje ćelija, 
formule te SXSSF Streaming API za velike datoteke.

\section{Word dokumenti}
Rukovanje .doc (HWPF) i .docx (XWPF) formatima, paragrafi, tablice i stilovi.

\section{PowerPoint prezentacije}
Rukovanje .ppt (HSLF) i .pptx (XSLF) formatima, slajdovi i oblici.

\chapter{Praktični primjeri}

\section{Generiranje Excel izvještaja}
Primjer automatskog generiranja tablice i grafa iz podataka.

\section{Obrada Word dokumenta}
Primjer čitanja i izmjene postojećeg .docx dokumenta.

\chapter{Zaključak}
Sažetak ključnih spoznaja i smjernice za primjenu Apache POI-a.

\makebackmatter

\appendices

\chapter{Prilog 1}

\end{document}